\section{Analysis and Lessons Learned}
\subsection{Different improvements require different evidence}
The model comparison establishes improved compliance and authored-target matching on identical inputs. The fresh semantic experiment tests incremental generation and commitment on a different, frozen corpus. The overload experiment isolates queue freshness with controlled payloads. These results should be presented together without merging their denominators or attributing every improvement to one mechanism. In particular, changing checkpoints is not an adaptive-policy gain, and fresh queue output is not the same as the right joke emphasis.

\subsection{A proposal should not be discarded solely because speech continues}
The exploratory readiness failure exposes a useful systems distinction. A phrase may already be represented correctly in a generated record while the model declines to authorize display. Caching permits additional evidence to arrive without paying for another proposal or losing the earlier work. Removing the final flag also avoids explicitly telling the proposal generator that the sentence is incomplete. The new paired baselines share that prompt, so their within-holdout comparison measures policy behavior; it does not separately identify the effect of each prompt change.

Source stability is an observable constraint rather than a semantic guarantee. A candidate can survive recognition revisions yet remain a poor emphasis, or become inappropriate when a later clause changes the interpretation. The ambiguity controls and pre-safe counts expose that risk. A bounded number of commits and a cooldown limit visible changes, but they can also prevent a later better proposal. The appropriate claim is a measured latency--coverage--workload tradeoff rather than universal dominance.

\subsection{Freshness control is useful when computation queues}
At light load, latest-only scheduling has little queued work to remove. The controlled overload ablation demonstrates where it matters: waiting jobs can become obsolete while still consuming FIFO service. Replacing them improves fresh-supported release coverage even though fewer total jobs complete. Deadline guards add a different benefit by avoiding or withholding work that cannot satisfy an age budget. Their retained fresh-output count is more informative than zero stale output alone.

The failure settings are equally useful for design. A conservative estimate greater than every deadline causes complete omission. A production scheduler should maintain request-dependent cost estimates and uncertainty, provide calibrated fallbacks, and expose coverage rather than treating suppression as successful service. These needs follow from the measured and replayed behavior rather than a claim about untested hardware.

\subsection{Progress beyond the checkpoint}
The status report's shared-clock concern has become a concrete source-level integration boundary and a new uncertain-input experiment. Current work adds measured component records, a reusable online policy interface, frozen development and holdout procedures, and independent causal checks. Positive findings support specific implementation choices while the initial failures explain the revision. No planned participant session is retroactively recorded as completed, and software evidence is kept separate from headset acceptance.
